\documentclass{article}
\usepackage{pstricks}
\usepackage{color}
\usepackage{pst-plot}
\title{PSTricks}
\author{Aditya Kumar Soni}
\date{\today}
% Since These compilers do not support pdf latex compiling, and pstricks uses a
post script compiler. so in order to rum the code, first run it in "LATEX" then
run using "DVI -> PS" and then "PS -> PDF" further using view pdf we can view the
document.
\begin{document}
\maketitle
\section{Creating Circle Graphics}
\vspace{4 cm}
\psgrid(4,4)
\pscircle(1,2){2}
\vspace{2cm}
\begin{pspicture}(0,0)(5,5)
\pscircle(3,3){2}
\psgrid(4,4)
\end{pspicture}
%for centering
\begin{center}
\psdots(1,0)(4,0)(1,4)
\end{center}
\newpage
\textbf{Exercise 1: Drawing a circle. Create a PSTricks figure that draws a
circle with centre at (2,3), and radius 1.5. Also plot the center point \newline
Answer:}
\begin{pspicture}(0,0)(5,5)
\pscircle(2,3){1.5}
\psgrid(5,5)
\psdots(2,3)
\end{pspicture}
\begin{center}
\psdots(2,3)
\psgrid(5,5)
\end{center}
\section{Plotting lines}
\begin{center}
\begin{pspicture}
%\psgrid(11,11)
\psdots(0,0)(7,0.5)(5,5)
\psline[linestyle=dashed, dash= 2pt 5pt, linewidth= 1pt, arrows=->](0,0)(7,0.5)
\psline[linestyle=dotted, dash= 10pt 5pt, linewidth= 1pt, arrows=->](0,0)(5,5)
\end{pspicture}
\end{center}
\newpage
\section{Plotting triangle and circles}
\begin{pspicture}
\psline[linecolor = black, fillstyle = solid, fillcolor = red](1,1)(5,1)(1,4)
(1,1)
\pscircle[linestyle = dotted](3,2.5){2.5}
\pscircle[linestyle = solid] (3,2.5){2.5}
\pscircle[fillstyle = solid, fillcolor = lightgray](2,2){1}
\psdots(2,2)
\end{pspicture}
\section{Plotting triangle and circles}
\begin{pspicture}
\psgrid(6,6)
%\psline(1,1)(5,1)(1,4)(1,1)
\pscircle[linestyle = dotted, linewidth = 3pt](3,2.5){2.5}
\pscircle[linecolor = red, fillstyle = hlines, fillcolor = lightgray](2,2){1}
\psdots(2,2)
\end{pspicture}
\newpage
\section{PRACTICAL Day 2 {Online class}}
\subsection{Compass Construction of an equilateral Triangle}
%\psgrid(5,5)
\begin{pspicture}
%\psgrid(10,10)
\psset{unit = 1.5} %this command changes the default unit of grid from 1 cm to
1.5 cm
%\pscircle[linewidth =2pt,linecolor =green](1,1){2}
\psarc[linewidth =3pt, linecolor =red](1,1){2}{0}{75}
%This command draws the arc of circle of radius 2 cm centered at point (1,1)
going from reference angle of 0 to 75 degree.
\psarc[linewidth =2pt](3,1){2}{110}{180}
\psline(1,1)(3,1)(2,2.732)(1,1)
\end{pspicture}
\subsection{Circle with shaded sector}
\vspace{4cm}
\psarc(2,2){1.5}{0}{60}
\pswedge[fillstyle=solid, fillcolor=red](2,2){1.5}{0}{60}
\begin{center}
\begin{pspicture}
%\psgrid{7,7}
\psset{unit=2}
\pscircle(2,2){1.5}
\pswedge[fillstyle=solid, fillcolor=red](2,2){1,5}{0}{60}
\put(2.75,1.7){$r$} 
%"\put" command tells where these graphical elements are placed
\put(2.3,2.1){$\theta$}
\put(1.75,1.70){0}
\end{pspicture}
\end{center}
\subsection{Plotting Functions}
% In order to plot functions using pstricks, use "\usepackage{pst-plot}" in the
preamble of document.
%\psplot[options]{xmin}{xmax}{xpression}
Plot the function $y=x$ from $x_{min}=-0.5$ to $x_{max}=4$
\newline
\begin{pspicture}(0,0)(5,5)
\psset{xunit=1cm,yunit=1cm}
\psaxes{->}(0,0)(-1,-1)(4.5,4.5)
%(0,0) is the origin, (-1,1) is lower left limit of axis, (4.5,4.5) is upper
right limit of axes
\psplot[linecolor=red]{-0.5}{4}{x}
\put(4.1,3.8){$y=x$}
\end{pspicture}
\newpage
Plot the function $y=\sin\frac{1}{x}$ from $x_{min}=0.025$ to $x_{max}=2$
\newline
\begin{pspicture}(-0.025,-4.25)(7.5,4.25)
\psset{xunit=3cm,yunit=3cm}
\psaxes{->}(0,0)(0,-1.25)(2.25,1.25)
\psplot[linecolor=green, plotpoints=2500]{0.025}{2}{1 x div RadtoDeg sin}
%this is to write sin(1/x), here div means division. This is written in
postscript (Reverse Polish Notation}
\put(7,-0.5){$x$}
\put(-0.5,4){$y$}
\put(3,3){$y=\sin\frac{1}{x}$}
\end{pspicture}
\begin{pspicture}(-0.025,-4.25)(7.5,4.25)
\psset{xunit=3cm,yunit=3cm}
\psaxes{->}(0,0)(0,-1.25)(2.25,1.25)
\psplot[linecolor=green, plotpoints=50]{0.025}{2}{1 x div RadtoDeg sin}
%this is to write sin(1/x), here div means division. This is written in
postscript (Reverse Polish Notation}
\put(7,-0.5){$x$}
\put(-0.5,4){$y$}
\put(3,3){$y=\sin\frac{1}{x}$}
\end{pspicture}
\end{document}